%% toc: P15 162 162.1,162.3,162.4
\section{Adaptedness, sign reversal and stopped information}
\label{sec:annex-p15-w02c-p15-019}

Throughout, \((\Omega,\mathcal F,P)\) is a probability space with a
filtration \((\mathcal F_n)_{n\geq0}\). The rules used for conditional
expectations are existence and uniqueness of versions
(\cref{thm:p15-conditional-existence}), linearity, order and the
known-information rule (\cref{thm:p15-conditional-basic}), the tower
identities (\cref{thm:p15-conditional-tower}) and the multiplier rule
(\cref{thm:p15-conditional-pullout}). A function measurable for the trivial field
\(\{\varnothing,\Omega\}\) is constant, because each of its sublevel
sets is empty or the whole space. For an integrable \(V\), the constant
\(EV\) is therefore the only version of
\(E(V\mid\{\varnothing,\Omega\})\): a version is constant, and its
integral over \(\Omega\) must be \(EV\).

\begin{annexitem}{annex:p15-w02c-p15-019-adaptedness-needed}{thm:p15-w02c-p15-019-adaptedness-needed}{Identity versus inequality}
\begin{proof}
\emph{Item~1.} Fix \(n\) and let \(Y_n\) be a version of
\(E(X_{n+1}\mid\mathcal F_n)\). It is \(\mathcal F_n\)-measurable and
integrable by \cref{thm:p15-conditional-existence}, and condition (iii)
says \(X_n=Y_n\) almost surely. Suppose now that \(\mathcal F_n\)
contains every event of \(\mathcal F\) of probability zero. For a Borel
set \(B\), both \(\{X_n\in B\}\) and \(\{Y_n\in B\}\) lie in
\(\mathcal F\), so their symmetric difference \(D\) lies in
\(\mathcal F\). It is contained in \(\{X_n\neq Y_n\}\), hence
\(P(D)=0\) and \(D\in\mathcal F_n\). Then
\(\{X_n\in B\}=\{Y_n\in B\}\mathbin{\triangle}D\in\mathcal F_n\), so
\(X_n\) is \(\mathcal F_n\)-measurable.

\emph{Item~2.} Since \(0\leq U\leq1\), \(|X_0|\leq3\), and every
\(X_n\) with \(n\geq1\) is the constant \(-1\); all are integrable,
which is condition (i). Every \(\mathcal F_n\) is trivial, so
\(E(X_{n+1}\mid\mathcal F_n)=EX_{n+1}=-1\) for every \(n\geq0\). For
\(n=0\) this gives
\[
       E(X_1\mid\mathcal F_0)=-1\geq-2-U=X_0
\]
at every point, because \(U\geq0\). For \(n\geq1\) it gives
\(E(X_{n+1}\mid\mathcal F_n)=-1=X_n\). Thus the submartingale form of
(iii) holds.

The \(\mathcal F_0\)-measurable variables are the constants. For every
real \(c\), \(P(X_0=c)=P(U=-2-c)=0\), since Lebesgue measure gives
each point measure zero. Hence \(X_0\) is not \(\mathcal F_0\)-measurable
and agrees almost surely with no \(\mathcal F_0\)-measurable variable;
in particular \(X\) is not adapted.

Finally \(EU=\int_0^1\omega\,d\omega=\tfrac12\), so
\(E(X_0\mid\mathcal F_0)=EX_0=-\tfrac52\) at every point, this being the
only version. The inequality \(-\tfrac52<-2-U\) holds exactly when
\(U<\tfrac12\), an event of probability \(\tfrac12\). The submartingale
form of \cref{thm:p15-martingale-conditioning} with \(m=n=0\) asserts
\(E(X_0\mid\mathcal F_0)\geq X_0\) almost surely, which therefore fails.
\end{proof}
\end{annexitem}

\begin{annexitem}{annex:p15-w02c-p15-019-multistep-definition}{thm:p15-w02c-p15-019-multistep-definition}{Gaps}
\begin{proof}
Conditions (i) and (ii) of \cref{def:p15-w02c-p15-019-adapted-classes}
are assumed, so only (iii) is in question.

Suppose \(X\) is a supermartingale. We prove
\(E(X_{m+k}\mid\mathcal F_m)\leq X_m\) almost surely by induction on
\(k\geq1\), for every \(m\). For \(k=1\) this is (iii). Let \(k\geq2\)
and assume the claim for \(k-1\) and every \(m\). Since
\(\mathcal F_m\subseteq\mathcal F_{m+k-1}\), the tower identity of
\cref{thm:p15-conditional-tower} gives
\[
 E(X_{m+k}\mid\mathcal F_m)
   =E\bigl(E(X_{m+k}\mid\mathcal F_{m+k-1})\bigm|\mathcal F_m\bigr)
\]
almost surely. By (iii) at time \(m+k-1\),
\(E(X_{m+k}\mid\mathcal F_{m+k-1})\leq X_{m+k-1}\) almost surely; both
sides are integrable, so the order rule of
\cref{thm:p15-conditional-basic} gives
\[
 E\bigl(E(X_{m+k}\mid\mathcal F_{m+k-1})\bigm|\mathcal F_m\bigr)
   \leq E(X_{m+k-1}\mid\mathcal F_m)\leq X_m ,
\]
the last step by the induction hypothesis with gap \(k-1\). This proves
the multi-step inequality for every \(n=m+k>m\).

Conversely, if the multi-step inequality holds for all \(n>m\), the
choice \(n=m+1\) is (iii) at time \(m\), so \(X\) is a supermartingale.

For submartingales, \(X\) is a submartingale if and only if \(-X\) is a
supermartingale (\cref{thm:p15-w02c-p15-019-sign-duality}), and by
linearity \(E(-X_n\mid\mathcal F_m)\leq-X_m\) almost surely exactly when
\(E(X_n\mid\mathcal F_m)\geq X_m\) almost surely; the supermartingale
equivalence for \(-X\) therefore gives the submartingale equivalence
for \(X\). For martingales, \(X\) is a martingale if and only if it is
both a submartingale and a supermartingale (the same lemma), and the
multi-step identity holds if and only if both multi-step inequalities
hold, the exceptional null sets being combined over the countably many
pairs \((m,n)\).
\end{proof}
\end{annexitem}

\begin{annexitem}{annex:p15-w02c-p15-019-sign-duality}{thm:p15-w02c-p15-019-sign-duality}{Negation preserves conditions (i) and (ii) and reverses (iii)}
\begin{proof}
Conditions (i) and (ii) of \cref{def:p15-w02c-p15-019-adapted-classes}
hold for \(X\) exactly when they hold for \(-X\): \(|-X_n|=|X_n|\), and
\(\{-X_n\in B\}=\{X_n\in-B\}\) for every Borel set \(B\), where
\(-B=\{-x:x\in B\}\) is Borel.

By linearity, \(E(-X_{n+1}\mid\mathcal F_n)=-E(X_{n+1}\mid\mathcal F_n)\)
almost surely. Hence, outside a null set,
\[
 E(X_{n+1}\mid\mathcal F_n)\geq X_n
 \quad\Longleftrightarrow\quad
 E(-X_{n+1}\mid\mathcal F_n)\leq -X_n .
\]
So the submartingale form of (iii) for \(X\) at time \(n\) holds almost
surely exactly when the supermartingale form holds for \(-X\). Together
with the first paragraph this proves the first equivalence.

If \(X\) is a martingale, then at each \(n\) the identity
\(E(X_{n+1}\mid\mathcal F_n)=X_n\) almost surely implies both
inequalities almost surely, so \(X\) is a submartingale and a
supermartingale. Conversely, if at time \(n\) the lower inequality holds
outside a null set \(N_1\) and the upper inequality outside a null set
\(N_2\), equality holds outside the null set \(N_1\cup N_2\). Conditions
(i) and (ii) are common to the three classes, so \(X\) is a martingale.
\end{proof}
\end{annexitem}

\begin{annexitem}{annex:p15-w02c-p15-019-event-comparison}{thm:p15-w02c-p15-019-event-comparison}{Test events where the conditional expectation exceeds the bound}
\begin{proof}
\emph{Item~1.} Let \(W\) be a version of \(E(Y\mid\mathcal G)\); it is
integrable and \(\mathcal G\)-measurable, and
\(\int_AW\,dP=\int_AY\,dP\) for every \(A\in\mathcal G\). For an integer
\(r\geq1\) put
\[
                    A_r=\{W-Z>1/r\}.
\]
Since \(W\) and \(Z\) are \(\mathcal G\)-measurable, \(A_r\in\mathcal G\).
The hypothesis and the defining identity give
\[
 \int_{A_r}(W-Z)\,dP=\int_{A_r}Y\,dP-\int_{A_r}Z\,dP\leq0,
\]
where the subtraction is legitimate because \(W\), \(Y\) and \(Z\) are
integrable. On \(A_r\) the integrand exceeds \(1/r\), so the same
integral is at least \(P(A_r)/r\). Hence \(P(A_r)=0\) for every \(r\).
The event \(\{W>Z\}\) is the union of the \(A_r\), \(r\geq1\), and is
therefore null. Thus \(W\leq Z\) almost surely, and the conclusion does
not depend on the version chosen, since two versions agree almost surely.

\emph{Item~2.} The variables \(-Y\) and \(-Z\) satisfy the hypotheses
of item~1, because \(\int_A(-Y)\,dP\leq\int_A(-Z)\,dP\) for every
\(A\in\mathcal G\). Item~1 gives \(E(-Y\mid\mathcal G)\leq-Z\) almost
surely, and linearity turns this into \(E(Y\mid\mathcal G)\geq Z\).

\emph{Item~3.} Equality for every \(A\) gives both hypotheses, hence both
inequalities almost surely, and so equality almost surely. Equivalently,
\(Z\) then satisfies the definition of a version of
\(E(Y\mid\mathcal G)\) in \cref{def:p15-general-conditional-expectation}.
\end{proof}
\end{annexitem}

\begin{annexitem}{annex:p15-w02c-p15-019-reduced-filtration}{thm:p15-w02c-p15-019-reduced-filtration}{The submartingale and martingale chains}
\begin{proof}
Condition (i) does not involve the filtration, and condition (ii) for
\((\mathcal G_n)\) is a hypothesis. Only (iii) remains. Fix \(n\).
Since \(\mathcal G_n\subseteq\mathcal F_n\), the tower identity of
\cref{thm:p15-conditional-tower} gives
\[
 E(X_{n+1}\mid\mathcal G_n)
     =E\bigl(E(X_{n+1}\mid\mathcal F_n)\bigm|\mathcal G_n\bigr)
\]
almost surely.

\emph{Submartingale.} By hypothesis
\(E(X_{n+1}\mid\mathcal F_n)\geq X_n\) almost surely, and both sides are
integrable. The order rule of \cref{thm:p15-conditional-basic} gives
\[
 E\bigl(E(X_{n+1}\mid\mathcal F_n)\bigm|\mathcal G_n\bigr)
       \geq E(X_n\mid\mathcal G_n)
\]
almost surely. Because \(X_n\) is \(\mathcal G_n\)-measurable and
integrable, the known-information rule gives
\(E(X_n\mid\mathcal G_n)=X_n\). Combining the three displays,
\(E(X_{n+1}\mid\mathcal G_n)\geq X_n\) almost surely.

\emph{Martingale.} Now \(E(X_{n+1}\mid\mathcal F_n)=X_n\) almost surely.
Two integrable variables that agree almost surely have the same
conditional expectation, by the order rule applied in both directions.
Hence
\[
 E\bigl(E(X_{n+1}\mid\mathcal F_n)\bigm|\mathcal G_n\bigr)
       =E(X_n\mid\mathcal G_n)=X_n ,
\]
and so \(E(X_{n+1}\mid\mathcal G_n)=X_n\) almost surely.

\emph{The same cases by sign reversal.} If \(X\) is a submartingale for
\((\mathcal F_n)\), then \(-X\) is a supermartingale for
\((\mathcal F_n)\) by \cref{thm:p15-w02c-p15-019-sign-duality}, and
\(-X\) is adapted to \((\mathcal G_n)\). The supermartingale case of
\cref{thm:p15-martingale-conditioning} makes \(-X\) a supermartingale
for \((\mathcal G_n)\), and the lemma again makes \(X\) a submartingale
for \((\mathcal G_n)\). A martingale for \((\mathcal F_n)\) is a
submartingale and a supermartingale for \((\mathcal F_n)\); each property
passes to \((\mathcal G_n)\), and the lemma makes \(X\) a martingale
for \((\mathcal G_n)\).

\emph{Natural filtration.} For \(i\leq n\), \(X_i\) is
\(\mathcal F_i\)-measurable and \(\mathcal F_i\subseteq\mathcal F_n\), so
\(\sigma(X_0,\ldots,X_n)\subseteq\mathcal F_n\). These fields increase
with \(n\), and \(X\) is adapted to them by construction, so the
theorem applies to them in each class.
\end{proof}
\end{annexitem}

\begin{annexitem}{annex:p15-w02c-p15-019-transform-cases}{thm:p15-w02c-p15-019-transform-cases}{Measurability, integrability and the nonnegative-stake inequality}
\begin{proof}
\emph{Item~1.} For \(i\leq n\), \(H_i\) is
\(\mathcal F_{i-1}\)-measurable and \(X_i,X_{i-1}\) are
\(\mathcal F_i\)-measurable by adaptedness; all three are
\(\mathcal F_n\)-measurable, and so is the finite sum
\((H\cdot X)_n=\sum_{i=1}^nH_i(X_i-X_{i-1})\). As in the proof of the
transform theorem, an almost-everywhere bound may be replaced by an
everywhere bound after clipping \(H_i\) to \([-K_i,K_i]\), which keeps
its measurability and changes the transform only on a null set. Then
\[
 E|(H\cdot X)_n|\leq\sum_{i=1}^nK_i\bigl(E|X_i|+E|X_{i-1}|\bigr)<\infty .
\]

\emph{Item~2.} In both constructions \(H_1=1\) is a constant, so it is
\(\mathcal F_0\)-measurable, and \((H\cdot X)_1=X_1-X_0=X_1\). In the
first, \(X_1=U\) is integrable but not constant, so it is not
measurable for the trivial field \(\mathcal F_1\); here \(X\) is not
adapted. In the second, \(X_1\) is a Borel function of \(U\), hence
\(\sigma(U)\)-measurable, and
\(E|X_1|=\int_0^1\omega^{-1}\,d\omega=\infty\); here \(X\) is adapted
but \(X_1\) is not integrable. In each case the corresponding
conclusion of item~1 fails at time \(1\).

\emph{Item~3.} Write \(G_n=(H\cdot X)_n\). By item~1, \(G\) is adapted
and integrable. Fix \(n\). Since \(G_{n+1}=G_n+H_{n+1}(X_{n+1}-X_n)\),
linearity gives
\[
 E(G_{n+1}\mid\mathcal F_n)
   =E(G_n\mid\mathcal F_n)+E\bigl(H_{n+1}(X_{n+1}-X_n)\bigm|\mathcal F_n\bigr).
\]
The first term is \(G_n\), because \(G_n\) is \(\mathcal F_n\)-measurable
and integrable. In the second, \(H_{n+1}\) is
\(\mathcal F_n\)-measurable and bounded, and \(X_{n+1}-X_n\) and its
product with \(H_{n+1}\) are integrable, so the multiplier rule of
\cref{thm:p15-conditional-pullout} and the known-information rule for
\(X_n\) give
\[
 E\bigl(H_{n+1}(X_{n+1}-X_n)\bigm|\mathcal F_n\bigr)
   =H_{n+1}\bigl(E(X_{n+1}\mid\mathcal F_n)-X_n\bigr).
\]
For a submartingale the bracket is nonnegative almost surely, and
\(H_{n+1}\geq0\), so the product is nonnegative almost surely. Hence
\(E(G_{n+1}\mid\mathcal F_n)\geq G_n\) almost surely, and \(G\) is a
submartingale.

The same conclusion follows by sign reversal. By
\cref{thm:p15-w02c-p15-019-sign-duality}, \(-X\) is a supermartingale.
The transform theorem \cref{thm:p15-predictable-transform}, with
\(H_i\geq0\), makes \(H\cdot(-X)\) a supermartingale. Since
\(H\cdot(-X)=-(H\cdot X)\) term by term, the lemma makes \(H\cdot X\) a
submartingale.
\end{proof}
\end{annexitem}

\begin{annexitem}{annex:p15-w02c-p15-019-sampling-cases}{thm:p15-w02c-p15-019-sampling-cases}{The event identity, each nonnegative term, and the comparison lemma}
\begin{proof}
By \cref{thm:p15-stopping-calculus}, \(X_S\) and \(X_T\) are integrable,
and \(X_S\) is \(\mathcal F_S\)-measurable; both stopping times are
bounded by \(N\) at every point. Fix \(A\in\mathcal F_S\). At each
point, \(X_T-X_S\) is the sum of the increments \(\Delta X_i\) over the
integers \(i\) with \(S<i\leq T\), all of which satisfy \(i\leq N\).
Hence
\[
 \mathbf1_A(X_T-X_S)=\sum_{i=1}^NC_i\,\Delta X_i,
 \qquad C_i=\mathbf1_{A\cap\{S\leq i-1\}\cap\{T\geq i\}} .
\]
Each \(C_i\) is \(\mathcal F_{i-1}\)-measurable:
\(A\cap\{S\leq i-1\}\in\mathcal F_{i-1}\) by the definition of
\(\mathcal F_S\), and \(\{T\geq i\}=\{T\leq i-1\}^c\in\mathcal F_{i-1}\).
Each term \(C_i\Delta X_i\) is integrable because \(0\leq C_i\leq1\).

Fix \(i\). The defining identity of \(E(X_i\mid\mathcal F_{i-1})\),
extended from indicators of \(\mathcal F_{i-1}\)-events to the bounded
\(\mathcal F_{i-1}\)-measurable \(C_i\) by the multiplier rule, gives
\(E(C_iX_i)=E\bigl(C_iE(X_i\mid\mathcal F_{i-1})\bigr)\). Because
\(X_{i-1}\) is \(\mathcal F_{i-1}\)-measurable,
\[
 E(C_i\Delta X_i)
   =E\Bigl[C_i\bigl(E(X_i\mid\mathcal F_{i-1})-X_{i-1}\bigr)\Bigr].
\]
For a submartingale the bracket is nonnegative almost surely and
\(C_i\geq0\), so \(E(C_i\Delta X_i)\geq0\). Summing over
\(1\leq i\leq N\),
\[
 \int_AX_T\,dP-\int_AX_S\,dP=E\bigl[\mathbf1_A(X_T-X_S)\bigr]\geq0 .
\]
This holds for every \(A\in\mathcal F_S\). Item~2 of
\cref{thm:p15-w02c-p15-019-event-comparison}, with
\(\mathcal G=\mathcal F_S\), \(Y=X_T\) and \(Z=X_S\), gives
\(E(X_T\mid\mathcal F_S)\geq X_S\) almost surely.

For a martingale each bracket is zero almost surely, so every
\(E(C_i\Delta X_i)\) is zero and \(\int_AX_T\,dP=\int_AX_S\,dP\) for every
\(A\in\mathcal F_S\); item~3 of the comparison lemma gives
\(E(X_T\mid\mathcal F_S)=X_S\) almost surely.

For a supermartingale each bracket is nonpositive, the event integrals
satisfy \(\int_AX_T\,dP\leq\int_AX_S\,dP\), and item~1 of the
comparison lemma gives \(E(X_T\mid\mathcal F_S)\leq X_S\); this is the
case written in \cref{annex:p15.bounded-optional-sampling}, and it also
yields the submartingale case for \(-X\) through
\cref{thm:p15-w02c-p15-019-sign-duality}.
\end{proof}
\end{annexitem}

\begin{annexitem}{annex:p15-w02c-p15-019-stopped}{thm:p15-w02c-p15-019-stopped-field-forms}{Differences and unions}
\begin{proof}
Let \(A\in\mathcal F\). Suppose \(A\cap\{T\leq n\}\in\mathcal F_n\) for
every \(n\). For \(n=0\), \(A\cap\{T=0\}=A\cap\{T\leq0\}\in\mathcal F_0\).
For \(n\geq1\),
\[
 A\cap\{T=n\}=\bigl(A\cap\{T\leq n\}\bigr)\setminus
                 \bigl(A\cap\{T\leq n-1\}\bigr),
\]
where the first event is in \(\mathcal F_n\) and the second is in
\(\mathcal F_{n-1}\subseteq\mathcal F_n\). Conversely, if
\(A\cap\{T=j\}\in\mathcal F_j\) for every finite \(j\), then
\[
 A\cap\{T\leq n\}=\bigcup_{j=0}^nA\cap\{T=j\}\in\mathcal F_n,
\]
since \(\mathcal F_j\subseteq\mathcal F_n\) for \(j\leq n\). The two
collections are therefore equal.

For finite \(k\) and \(n\), \(\{T=k\}\cap\{T=n\}\) is \(\{T=k\}\) when
\(n=k\) and empty otherwise; in both cases it lies in \(\mathcal F_n\),
because \(\{T=k\}\in\mathcal F_k\) by \cref{thm:p15-stopping-calculus}.
So \(\{T=k\}\in\mathcal F_T\). The event \(\{T\leq k\}\) is the finite
union of \(\{T=j\}\), \(j\leq k\), and \(\{T=\infty\}\) is the complement
of the countable union of all \(\{T=j\}\). Since \(\mathcal F_T\) is a
sigma-algebra by \cref{thm:p15-stopping-calculus}, both lie in
\(\mathcal F_T\). The sets \(\{k\}\), \(k\) finite, and \(\{\infty\}\)
generate all subsets of \(\{0,1,\ldots\}\cup\{\infty\}\) under countable
unions, so every inverse image under \(T\) lies in \(\mathcal F_T\).
\end{proof}
\end{annexitem}

\begin{annexitem}{annex:p15-w02c-p15-019-crossing}{thm:p15-w02c-p15-019-crossing-needs-adapted}{A nontrivial first purchase}
\begin{proof}
Since \(X_m=U\leq1<2=b\) at every point, no time ever satisfies
\(X_m\geq b\). With \(a=\tfrac12\), the definition gives \(T_1=0\) on
\(\{U\leq\tfrac12\}\) and \(T_1=\infty\) on \(\{U>\tfrac12\}\), because
\(X_m=U\) for every \(m\). In either case \(T_2=\infty\), and every later
alternating time is infinite. Hence
\[
                     \{T_1\leq0\}=\{U\leq\tfrac12\},
\]
an event of probability \(\tfrac12\), which is neither empty nor the
whole space and so is not in the trivial field \(\mathcal F_0\). Thus
\(T_1\) is not a stopping time. For \(m=1\), the only interval that can
contain \(1\) is the first, and \(T_1<1\leq T_2\) holds exactly when
\(T_1=0\). So \(H_1=\mathbf1_{\{U\leq1/2\}}\), which is not constant and
hence not \(\mathcal F_0\)-measurable.
\end{proof}
\end{annexitem}

\begin{annexitem}{annex:p15-w02c-p15-019-mean-value}{thm:p15-w02c-p15-019-mean-value}{A ball average in the radius}
\begin{proof}
Fix \(x\in\mathbb R^d\). Write \(\partial_i\) for the partial derivative
in the \(i\)th coordinate and
\(g=\partial_1\partial_1f+\cdots+\partial_d\partial_df\), so that
\(g\leq0\) everywhere and \(g\) is continuous. Let
\(\bar B=\{z:|z|\leq1\}\) be the closed unit ball and
\(\beta=|\bar B|\in(0,\infty)\) its volume. For \(r\geq0\) put
\[
            \psi(r)=\frac1\beta\int_{\bar B}f(x+rz)\,dz .
\]
\emph{Step 1: \(\psi\) is the ball average and \(\psi(0)=f(x)\).} For
\(r>0\) the map \(z\mapsto x+rz\) carries \(\bar B\) onto
\(B(x,r)\) and multiplies Lebesgue measure by \(r^d\), so
\(|B(x,r)|=r^d\beta\) and
\[
     \psi(r)=\frac1{r^d\beta}\int_{B(x,r)}f(y)\,dy
            =\frac1{|B(x,r)|}\int_{B(x,r)}f(y)\,dy .
\]
At \(r=0\) the integrand is the constant \(f(x)\), so \(\psi(0)=f(x)\).
It therefore suffices to show \(\psi(r)\leq\psi(0)\) for \(r>0\).

\emph{Step 2: \(\psi\) is continuous on \([0,\infty)\) and
differentiable on \((0,\infty)\).} Fix \(R>0\). The set
\(Q=\{x+sz:0\leq s\leq R+1,\ z\in\bar B\}\) is compact, so \(f\) and
its first partial derivatives are bounded on \(Q\), say by \(M\). If
\(r_k\to r\) in \([0,R]\), then \(f(x+r_kz)\to f(x+rz)\) for every
\(z\) by continuity, with \(|f(x+r_kz)|\leq M\); dominated convergence
on the finite-measure set \(\bar B\) gives \(\psi(r_k)\to\psi(r)\). For
the derivative, fix \(r\in(0,R)\) and \(0<|h|<\min(r,R-r)\). For each
\(z\) the function \(s\mapsto f(x+sz)\) has derivative
\(\sum_i z_i\,\partial_if(x+sz)\) by the chain rule, and the mean value
theorem gives
\[
   \Bigl|\frac{f(x+(r+h)z)-f(x+rz)}h\Bigr|\leq\sum_{i=1}^d|z_i|\,M\leq dM .
\]
The difference quotients converge pointwise to
\(\sum_iz_i\,\partial_if(x+rz)\), so dominated convergence gives
\[
   \psi'(r)=\frac1\beta\int_{\bar B}\sum_{i=1}^dz_i\,\partial_if(x+rz)\,dz .
\]

\emph{Step 3: \(\psi'(r)\leq0\) for \(r>0\).} For each \(i\) define on
\(\mathbb R^d\)
\[
             V_i(z)=\partial_if(x+rz)\,\bigl(1-|z|^2\bigr).
\]
Because \(\partial_if\) has continuous first partial derivatives, the
chain and product rules give the continuous partial derivative
\[
  \frac{\partial V_i}{\partial z_i}(z)
     =r\,\partial_i\partial_if(x+rz)\bigl(1-|z|^2\bigr)
       -2z_i\,\partial_if(x+rz).
\]
We show that its integral over \(\bar B\) is zero. If \(d=1\), then
\(\bar B=[-1,1]\) and the fundamental theorem of calculus gives
\(V_1(1)-V_1(-1)=0-0=0\). If \(d\geq2\), write \(z=(t,z')\) with \(t\)
the \(i\)th coordinate and \(z'\in\mathbb R^{d-1}\) the others. Then
\(\bar B=\{(t,z'):|z'|\leq1,\ |t|\leq a(z')\}\) with
\(a(z')=(1-|z'|^2)^{1/2}\). The integrand is continuous on the compact
set \(\bar B\), hence bounded and integrable, and Fubini's theorem gives
\[
 \int_{\bar B}\frac{\partial V_i}{\partial z_i}\,dz
   =\int_{\{|z'|\leq1\}}\Bigl(\int_{-a(z')}^{a(z')}
        \frac{\partial V_i}{\partial t}(t,z')\,dt\Bigr)dz' .
\]
For each \(z'\) the inner integral equals
\(V_i(a(z'),z')-V_i(-a(z'),z')\) by the fundamental theorem of
calculus. Both endpoints have \(|z|=1\), where \(V_i\) vanishes, so the
inner integral is zero, and so is the whole integral.

Summing over \(i\) and using the formula for
\(\partial V_i/\partial z_i\),
\[
 0=r\int_{\bar B}g(x+rz)\bigl(1-|z|^2\bigr)\,dz
   -2\int_{\bar B}\sum_{i=1}^dz_i\,\partial_if(x+rz)\,dz .
\]
On \(\bar B\), \(1-|z|^2\geq0\) and \(g(x+rz)\leq0\), and \(r>0\), so
the first integral is at most zero. Hence
\[
   \beta\,\psi'(r)=\int_{\bar B}\sum_{i=1}^dz_i\,\partial_if(x+rz)\,dz
      =\frac r2\int_{\bar B}g(x+rz)\bigl(1-|z|^2\bigr)\,dz\leq0 .
\]

\emph{Step 4: conclusion.} Let \(r>0\). By Step~2, \(\psi\) is continuous
on \([0,r]\) and differentiable on \((0,r)\), so the mean value theorem
gives some \(s\in(0,r)\) with
\[
          \psi(r)-\psi(0)=r\,\psi'(s)\leq0
\]
by Step~3. With Step~1,
\(|B(x,r)|^{-1}\int_{B(x,r)}f(y)\,dy=\psi(r)\leq\psi(0)=f(x)\).
\end{proof}
\end{annexitem}

\begin{annexitem}{annex:p15-w02c-p15-019-doubling-unfavorable}{thm:p15-w02c-p15-019-doubling-unfavorable}{Doubling stakes}
\begin{proof}
\emph{The stakes.} The doubling system is \(H_1=1\) and, for
\(n\geq2\),
\(H_n=2H_{n-1}\mathbf1_{\{\xi_{n-1}=-1\}}+\mathbf1_{\{\xi_{n-1}=1\}}\);
exactly one indicator equals \(1\) at each point because \(\xi_{n-1}\)
takes only the values \(\pm1\). By induction, \(H_n\) is
\(\mathcal F_{n-1}\)-measurable and \(1\leq H_n\leq2^{n-1}\): this holds
for \(n=1\) since \(H_1\) is constant, and if it holds for \(n-1\), then
\(H_{n-1}\) and \(\xi_{n-1}\) are \(\mathcal F_{n-1}\)-measurable and
\(H_n\) lies between \(1\) and \(\max(1,2\cdot2^{n-2})=2^{n-1}\).

\emph{The walk.} Each \(\xi_i\) is bounded with mean
\(\mu=p-(1-p)=2p-1\leq0\). The examples theorem
\cref{thm:p15-martingale-examples}, with \(c=0\), shows that
\(X_n=\xi_1+\cdots+\xi_n\) is a supermartingale for
\((\mathcal F_n)\), and a martingale when \(p=\tfrac12\). The theorem
assumes identically distributed increments, which holds here because
every \(\xi_n\) has the same law. \(X\) is adapted, since \(X_n\) is a
function of \(\xi_1,\ldots,\xi_n\).

\emph{The transform and the stopped transform.} The stakes are
predictable, nonnegative and bounded by the deterministic constants
\(2^{n-1}\), so by \cref{thm:p15-predictable-transform} the transform
\(H\cdot X\) is a supermartingale, and a martingale when
\(p=\tfrac12\); by item~1 of \cref{thm:p15-w02c-p15-019-transform-cases}
it is adapted and integrable. The time \(T\) is a stopping time, since
\(\{T\leq n\}=\bigcup_{m=1}^n\{\xi_m=1\}\in\mathcal F_n\). By
\cref{thm:p15-stopping-calculus}, stopping preserves each class, so
\(G_n=(H\cdot X)_{T\wedge n}\) is a supermartingale, and a martingale
when \(p=\tfrac12\). Both processes are \(0\) at time \(0\).

\emph{The expectations.} For a supermartingale \(Z\) with \(Z_0=0\),
taking expectations in \(E(Z_{n+1}\mid\mathcal F_n)\leq Z_n\) gives
\(EZ_{n+1}\leq EZ_n\), and by induction \(EZ_n\leq EZ_0=0\). For a
martingale the same steps give equality. Applying this to \(H\cdot X\)
and to \(G\) gives the stated expectations.

\emph{The almost-sure limit.} Since the \(\xi_m\) are independent,
\(P(T>n)=P(\xi_1=\cdots=\xi_n=-1)=(1-p)^n\), which tends to \(0\)
because \(p>0\); hence \(T<\infty\) almost surely. On \(\{T\leq n\}\),
\(G_n=(H\cdot X)_T\) is the gain after \(T-1\) losses and one win, which
is \(-(1+2+\cdots+2^{T-2})+2^{T-1}=1\), as in
\cref{thm:p15-martingale-boundary-calculations}. So for every \(\omega\)
with \(T(\omega)<\infty\), \(G_n(\omega)=1\) for all
\(n\geq T(\omega)\), and \(G_n\to1\) almost surely.
\end{proof}
\end{annexitem}

\begin{annexitem}{annex:p15-w02c-p15-019-no-loss}{thm:p15-w02c-p15-019-no-loss}{Nonnegative stakes}
\begin{proof}
A submartingale is adapted and integrable by
\cref{def:p15-w02c-p15-019-adapted-classes}, and \(K\) is predictable,
nonnegative and bounded by deterministic constants. Item~3 of
\cref{thm:p15-w02c-p15-019-transform-cases} makes \(K\cdot Y\) a
submartingale, and \((K\cdot Y)_0=0\) by
\cref{def:p15-predictable-transform}. Write \(Z=K\cdot Y\). Taking
expectations in \(E(Z_{n+1}\mid\mathcal F_n)\geq Z_n\) gives
\(EZ_{n+1}\geq EZ_n\), and by induction \(EZ_n\geq EZ_0=0\).
\end{proof}
\end{annexitem}
